%% Data identitas bersama seluruh dokumen di contoh/ (satu sumber di akar contoh/).
%% Dipakai semua contoh lewat identitas.tex masing-masing: proposal, laporan, maupun
%% repositori hanya menambah atau menimpa kunci yang khusus dokumen itu (judul, pembimbing
%% kedua, gelar). Cukup ubah berkas ini sekali untuk mengganti data pada semua contoh.
\identitas{
  nama               = {Harjito},
  nim                = {4314931284},
  prodi              = {Pendidikan Teknik Informatika},
  fakultas           = {Fakultas Teknik},
  kota               = {Semarang},
  tahun              = {2026},
  pembimbinga        = {Dr. Pembimbing Satu, M.Kom.},
  pembimbinganip     = {198001012005011001},
  ketua              = {Prof. Dr. Ketua Penguji, M.Pd.},
  ketuanip           = {196001011990031001},
  sekretaris         = {Dr. Sekretaris Penguji, M.Kom.},
  sekretarisnip      = {197002022000031002},
  penguji1           = {Dr. Penguji Satu, M.Pd.},
  penguji1nip        = {197503032005012003},
  penguji2           = {Dr. Penguji Dua, M.Kom.},
  penguji2nip        = {198004042010011004},
  hari               = {Senin},
  tanggal            = {12 Januari},
  tanggalpersetujuan = {10 Januari 2026},
  tanggalpernyataan  = {10 Januari 2026},
  tempatlahir        = {Semarang},
  tanggallahir       = {1 Januari 2004},
  alamat             = {Jalan Contoh Nomor 1, Sekaran, Gunungpati, Semarang},
  surel              = {nama.mahasiswa@students.unnes.ac.id}
}
